\begin{table}[t]
\centering
\tiny
\caption{Structural findings extracted from generated severity tables (table\_severity\_auc/table\_severity\_slope) and winner analyses.}
\label{tab:structural-findings}
\resizebox{\columnwidth}{!}{%
\begin{tabular}{p{0.24\linewidth}p{0.34\linewidth}p{0.31\linewidth}}
\toprule
Finding & Quantitative evidence & Deployment implication \\
\midrule
Metric-outage robustness is model-family specific & Salesforce/Borg missing-variable slope: DLinear 16.78 vs PatchTST-lite 8.40; Alibaba: 2.30 vs 1.12. & Patch context helps when entire metric channels disappear, but it should be weighed against latency. \\
Patch robustness is not universal & Noise slopes are close: Salesforce/Borg 0.78-0.99, Alibaba 0.05-0.07 for DLinear/RACE/PatchTST-lite. & A model that helps telemetry outage may not add much under point-wise sensor noise. \\
Clean, latency, and capacity objectives disagree & On Alibaba core learned scenarios, best MSE differs from best learned P95 in 4/5 cases and from best capacity proxy in 4/5 cases. & A clean leaderboard cannot decide between online alerting and proactive capacity planning. \\
Robustness has measurable deployment cost & On Salesforce/Borg metric outage, PatchTST-lite has about 3.6x DLinear P95 latency while roughly halving the degradation slope. & The benchmark should report Pareto choices, not only a single accuracy winner. \\
\bottomrule
\end{tabular}
}
\end{table}
